

\clearpage
\section{Objetivos}
La investigación se plantea con un propósito o finalidad, por lo que se requiere establecer lo que se pretende de manera particular. En este sentido, los objetivos deben ser concretos, medibles, alcanzables y realistas. Se redactan con un verbo que conduzca a la acción (en infinitivo) y guiarán el desarrollo del estudio, deberán ser coherentes entre sí \cite{ejemplo2020}. Según \cite{ejemplo2022} los objetivos constituyen la base para la estructuración de la metodología. Se establecen dos tipos de objetivos:
\subsection{Objetivo General}
Se formulará un solo objetivo general el cual debe mostrar una relación clara y coherente con la pregunta central de la investigación y evidenciará los logros que se pretende alcanzar con el proyecto de investigación \cite{ejemplo2021}.
\subsection{Objetivos Específicos}
Los objetivos específicos serán máximo dos y demostrarán los logros parciales del proceso de investigación de manera secuencial; permiten la operatividad del objetivo general y propician el cumplimiento de expectativas originadas en las variables o categorías de estudio. Los objetivos son las metas que deberán cumplirse para dar solución al problema y reflejarán los logros en el proceso secuencial de la investigación \cite{ejemplo2020}.
\subsection{Alcance de proyecto}
El alcance de un proyecto se refiere a la definición clara y detallada de lo que el proyecto va a entregar y lo que no. Es uno de los elementos clave de la planificación y responde a las preguntas:
    ¿Qué se va a hacer?
    ¿Qué se va a entregar como resultado final?
    ¿Qué actividades y tareas estarán incluidas?
    ¿Qué no estará incluido?
Ejemplo:

\subsubsection{Fase 1(Objetivo Específico 1)}
En esta subsección se debe describir cada una de las actividades que se van a cumplir en la primera fase del proyecto que se traduce en el objetivo específico 1
\paragraph{Actividad 1:} Por cada una de las fases se debe desarrollar 3 actividades como mínimo. En cada actividad se debe describir claramente las tareas que se van a realizar, así como las tecnologías que se va a utilizar para el desarrollo de la actividad y finalmente se debe nombrar cual es el entregable que se produce una vez culminada la actividad. Recuerde que una actividad se diferencia de una tarea por la magnitud y el tiempo desarrollo de la misma, la actividad toma más tiempo en ser desarrollada. 

\paragraph{Actividad 2:} Por cada una de las fases se debe desarrollar 3 actividades como mínimo. En cada actividad se debe describir claramente las tareas que se van a realizar, así como las tecnologías que se va a utilizar para el desarrollo de la actividad y finalmente se debe nombrar cual es el entregable que se produce una vez culminada la actividad. Recuerde que una actividad se diferencia de una tarea por la magnitud y el tiempo desarrollo de la misma, la actividad toma más tiempo en ser desarrollada.

\paragraph{Actividad 3:} Por cada una de las fases se debe desarrollar 3 actividades como mínimo. En cada actividad se debe describir claramente las tareas que se van a realizar, así como las tecnologías que se va a utilizar para el desarrollo de la actividad y finalmente se debe nombrar cual es el entregable que se produce una vez culminada la actividad. Recuerde que una actividad se diferencia de una tarea por la magnitud y el tiempo desarrollo de la misma, la actividad toma más tiempo en ser desarrollada.

\subsubsection{Fase 2(Objetivo Específico 2)}
En esta subsección se debe describir cada una de las actividades que se van a cumplir en la primera fase del proyecto que se traduce en el objetivo específico 2
\paragraph{Actividad 1:} Por cada una de las fases se debe desarrollar 3 actividades como mínimo. En cada actividad se debe describir claramente las tareas que se van a realizar, así como las tecnologías que se va a utilizar para el desarrollo de la actividad y finalmente se debe nombrar cual es el entregable que se produce una vez culminada la actividad. Recuerde que una actividad se diferencia de una tarea por la magnitud y el tiempo desarrollo de la misma, la actividad toma más tiempo en ser desarrollada. 

\paragraph{Actividad 2:} Por cada una de las fases se debe desarrollar 3 actividades como mínimo. En cada actividad se debe describir claramente las tareas que se van a realizar, así como las tecnologías que se va a utilizar para el desarrollo de la actividad y finalmente se debe nombrar cual es el entregable que se produce una vez culminada la actividad. Recuerde que una actividad se diferencia de una tarea por la magnitud y el tiempo desarrollo de la misma, la actividad toma más tiempo en ser desarrollada.

\paragraph{Actividad 3:} Por cada una de las fases se debe desarrollar 3 actividades como mínimo. En cada actividad se debe describir claramente las tareas que se van a realizar, así como las tecnologías que se va a utilizar para el desarrollo de la actividad y finalmente se debe nombrar cual es el entregable que se produce una vez culminada la actividad. Recuerde que una actividad se diferencia de una tarea por la magnitud y el tiempo desarrollo de la misma, la actividad toma más tiempo en ser desarrollada.

